\documentclass{article}
\usepackage[utf8]{inputenc}
\usepackage{amsmath}
\usepackage{amsfonts}
\usepackage{amssymb}
\usepackage{geometry}
\usepackage{tcolorbox}
\usepackage{listings}
\usepackage{xcolor}
\usepackage{hyperref}

\geometry{a4paper, margin=1in}

\definecolor{codegreen}{rgb}{0,0.6,0}
\definecolor{codegray}{rgb}{0.5,0.5,0.5}
\definecolor{codepurple}{rgb}{0.58,0,0.82}
\definecolor{backcolour}{rgb}{0.95,0.95,0.92}

\lstset{
    backgroundcolor=\color{backcolour},   
    commentstyle=\color{codegreen},
    keywordstyle=\color{magenta},
    numberstyle=\tiny\color{codegray},
    stringstyle=\color{codepurple},
    basicstyle=\ttfamily\small,
    breakatwhitespace=false,         
    breaklines=true,                 
    captionpos=b,                    
    keepspaces=true,                 
    numbers=left,                    
    numbersep=5pt,                  
    showspaces=false,                
    showstringspaces=false,
    showtabs=false,                  
    tabsize=2
}

\title{Module 04: Feature Engineering II - Target Encoding and the Power of Smoothing}
\author{Cybernauts 2026 Datathon Campaign}
\date{May 2026}

\begin{document}

\maketitle

\begin{abstract}
Target Encoding is arguably the most powerful encoding technique in a data scientist's arsenal, especially for high-cardinality features like \texttt{ZIP\_Code} or \texttt{Product\_ID}. However, it is a "double-edged sword." If used without caution, it introduces \textbf{Data Leakage} and causes your model to overfit instantly. This note covers how to implement Target Encoding safely using M-Estimate Smoothing.
\end{abstract}

\tableofcontents
\newpage

\section{The Core Concept: Encoding the Signal}
In earlier modules, we discussed One-Hot and Frequency encoding. \textbf{Target Encoding} goes a step further by using the information from the target variable ($y$) itself to encode the categories.

\textit{Logic:} You replace each text category with the \textbf{average target value} for that specific category.
\textit{Example:} If you are predicting "Loan Default" (0 or 1) and users in "Dhaka" default 15\% of the time, the category "Dhaka" is replaced by the number \textbf{0.15}.

\section{The Danger: Data Leakage and Overfitting}
The biggest risk of Target Encoding is \textbf{Overfitting}. 
\begin{tcolorbox}[colback=red!5!white,colframe=red!75!black,title=The Single-Row Trap]
Imagine a category called "Specific\_Device" that only appears \textbf{once} in your entire dataset. If that specific user happened to default (Target = 1), the encoded value for that category becomes exactly \textbf{1.0}. 
The model will see this 1.0 and think it has found a "perfect predictor," but in reality, it has just memorized a single row. This is a form of data leakage.
\end{tcolorbox}

\section{The Solution: M-Estimate Smoothing}
To fix the overfitting problem, we don't use the raw category average. Instead, we "smooth" it by blending it with the \textbf{Global Average} of the entire dataset.

\subsection{The Smoothing Formula}
The encoded value $S_i$ for a category $i$ is calculated as:
\begin{equation}
    S_i = \alpha_i \cdot \bar{y}_i + (1 - \alpha_i) \cdot \bar{y}_{\text{global}}
\end{equation}
Where:
\begin{itemize}
    \item $\bar{y}_i$: The average target value for the specific category.
    \item $\bar{y}_{\text{global}}$: The average target value for the \textbf{entire} training set.
    \item $\alpha_i$: A weight (between 0 and 1) based on the number of samples in that category.
\end{itemize}

\subsection{How Smoothing Behaves}
\begin{itemize}
    \item \textbf{High Count Categories:} If a category appears 10,000 times, $\alpha_i$ is close to 1. The encoded value stays very close to the specific category average ($\bar{y}_i$).
    \item \textbf{Low Count Categories (Rare):} If a category appears only 2 times, $\alpha_i$ is close to 0. The encoded value is "pulled" toward the global average ($\bar{y}_{\text{global}}$), preventing the model from overreacting to a tiny sample size.
\end{itemize}

\section{Python Implementation}
While you can calculate this manually, the \texttt{category\_encoders} library is the industry standard for safe implementation.

\begin{lstlisting}[language=Python]
from category_encoders import TargetEncoder

# Initialize encoder with smoothing parameter
te = TargetEncoder(smoothing=10)

# Fit and transform (IMPORTANT: Only fit on Training data to avoid leakage!)
X_train['City_Target'] = te.fit_transform(X_train['City'], y_train)
X_test['City_Target'] = te.transform(X_test['City'])
\end{lstlisting}

\section{Practice Problems}
\textbf{P1:} Why is it critical to only \texttt{fit} the TargetEncoder on the training set and not the full dataset? \\
\textit{Answer: To prevent \textbf{Data Leakage}. If you use the test set targets to calculate the encoding, the model "sees" the future during training and will give an artificially high accuracy that fails on real data.}

\textbf{P2:} If the Global Average Default Rate is 0.10, and a rare city "Khulna" appears only twice with a 0.50 default rate. Will a smoothed Target Encoder output 0.50? \\
\textit{Answer: No. Because it is a rare category, the encoder will pull the value down toward the global average (0.10), resulting in a value like 0.15 or 0.20 instead of 0.50.}

\end{document}
